\documentclass[11pt,a4paper]{article}
\usepackage{fontspec}
\usepackage{polyglossia}
\setdefaultlanguage{russian}
\setotherlanguage{english}
\setmainfont[Path=fonts/,BoldFont=PT_Serif-Web-Bold.ttf,ItalicFont=PT_Serif-Web-Italic.ttf,BoldItalicFont=PT_Serif-Web-BoldItalic.ttf]{PT_Serif-Web-Regular.ttf}
\setsansfont[Path=fonts/,BoldFont=PT_Sans-Web-Bold.ttf]{PT_Sans-Web-Regular.ttf}
\newfontfamily\cyrillicfontsf[Path=fonts/,BoldFont=PT_Sans-Web-Bold.ttf]{PT_Sans-Web-Regular.ttf}
\newfontfamily\narrowsf[Path=fonts/,BoldFont=PT_Sans-Narrow-Web-Bold.ttf]{PT_Sans-Narrow-Web-Regular.ttf}
\setmonofont[Path=fonts/,Scale=0.92]{PTMono.ttf}
\newfontfamily\tengw[Path=fonts/,Renderer=HarfBuzz]{AlcarinTengwar-Regular.otf}

\usepackage[a4paper,margin=22mm,top=24mm,bottom=26mm]{geometry}
\usepackage{xcolor}
\usepackage{tikz}
\usetikzlibrary{arrows.meta,positioning,calc,fit,backgrounds,patterns,decorations.pathreplacing,shapes.geometric,automata}
\usepackage{booktabs,longtable,array,multirow,tabularx}
\usepackage{caption}
\captionsetup{font=small,labelfont=bf,labelsep=period,skip=6pt}
\usepackage{enumitem}
\setlist{itemsep=2pt,topsep=4pt}
\usepackage{fancyhdr}
\usepackage[hidelinks]{hyperref}
\usepackage{amsmath}

% ---------- tengwar helpers ----------
\newlength\tgem \setlength\tgem{5mm}
\newcommand\tgu{0.001\tgem}
\newcommand\TG[1]{{\tengw\fontsize{\tgem}{\tgem}\selectfont\XeTeXglyph#1\relax}}
\newcommand\TGtext[1]{{\tengw\fontsize{\tgem}{\tgem}\selectfont #1}}
\newcommand\tgsmall{\tengw\fontsize{\tgem}{\tgem}\selectfont}
\newcommand\tgkeyL{\tengw\fontsize{7.5mm}{7.5mm}\selectfont}
\newcommand\tgkeyU{\tengw\fontsize{5.2mm}{5.2mm}\selectfont}
\newcommand\tgkuL{0.0075mm}
\newcommand\tgkuU{0.0052mm}
\newcommand\TGkL[1]{{\tengw\fontsize{7.5mm}{7.5mm}\selectfont\XeTeXglyph#1\relax}}
\newcommand\TGkU[1]{{\tengw\fontsize{5.2mm}{5.2mm}\selectfont\XeTeXglyph#1\relax}}
\newcommand\tg[1]{{\tengw\fontsize{1.35em}{1em}\selectfont #1}}  % inline tengwar in text

\definecolor{ink}{HTML}{1F2328}
\definecolor{axisred}{HTML}{B3261E}
\definecolor{axisblue}{HTML}{1F5FAD}
\definecolor{cellc}{HTML}{EEF1F5}
\definecolor{slugc}{HTML}{8A8F98}
\definecolor{dead}{HTML}{FBEFD5}
\definecolor{adv}{HTML}{E8EEF7}
\definecolor{fnk}{HTML}{E4E4E4}
\definecolor{accent}{HTML}{6B3FA0}

\tikzset{
  cellfill/.style={fill=cellc},
  cellline/.style={draw=slugc!70,thin},
  axisT/.style={draw=axisred,thick,dash pattern=on 6pt off 2pt on 1pt off 2pt},
  axisB/.style={draw=axisblue,thick,dashed},
  base/.style={draw=slugc,thin},
  slug/.style={draw=slugc,rounded corners=2pt,thin},
  dim/.style={draw=ink!70,{Latex[length=1.6mm]}-{Latex[length=1.6mm]},thin},
  lbl/.style={font=\sffamily\footnotesize,text=ink,inner sep=1pt},
  ghost/.style={opacity=0.22},
  key/.style={draw=ink!55,rounded corners=1.5pt,thin},
  keydead/.style={},keyadv/.style={},
  keyleg/.style={font=\ttfamily\tiny,text=ink!45,anchor=north west,inner sep=0pt},
  keyres/.style={font=\sffamily\tiny,text=ink!40},
  keyup/.style={text=accent},
  keylo/.style={text=ink},
  fnleg/.style={font=\sffamily\scriptsize,text=ink},
  blk/.style={draw=ink!70,rounded corners=2pt,fill=white,font=\sffamily\small,align=center,inner sep=4pt,minimum height=9mm},
  arr/.style={-{Latex[length=2mm]},thick,draw=ink!80},
  mech/.style={draw=ink!80,thick},
}

\pagestyle{fancy}\fancyhf{}
\renewcommand{\headrulewidth}{0.3pt}
\fancyhead[L]{\sffamily\small Тенгварская пишущая машинка ТПМ-1}
\fancyhead[R]{\sffamily\small Проектный документ · черновик 0.1}
\fancyfoot[C]{\sffamily\small\thepage}

\newcommand\code[1]{\texttt{#1}}
\newcommand\un{\,u}

\begin{document}

\begin{titlepage}
\vspace*{18mm}
{\sffamily\small ПРОЕКТНЫЙ ДОКУМЕНТ · ЧЕРНОВИК 0.1 · 10 НОЯБРЯ 2026}\par
\vspace{8mm}
{\sffamily\bfseries\fontsize{26}{30}\selectfont Тенгварская пишущая\\[2pt] машинка ТПМ-1}\par
\vspace{4mm}
{\large Осевые литеры, пропорциональная подача\\ с отложенной второй фазой}\par
\vspace{14mm}
\begin{center}
\setlength\tgem{15mm}
\begin{tikzpicture}[x=\tgu,y=\tgu]
\input{gen/fig_trace_q.tex}
\end{tikzpicture}
\end{center}
\vspace{6mm}
\noindent\begin{minipage}{0.86\textwidth}\small
\textbf{О чём документ.} Описание гипотетической электромеханической пишущей машинки для тенгвара в режимах с техтами (квенийском, синдаринском, общем). Главная идея: каждая литера отлита так, что ось её техты совпадает с осью рычага. Каретка ставит под удар именно эту ось, а не левый край клетки. Тогда один набор мёртвых техт годится для узких, широких и «смещённых» тенгв. Ширины и оси всех литер не придуманы, а измерены на цифровом шрифте Alcarin Tengwar и пересчитаны в единицы подачи. На титуле вверху — оттиск машинки, рассчитанный симулятором каретки; внизу — тот же текст, набранный шрифтом с полной OpenType-раскладкой техт.
\end{minipage}
\vfill
{\small User · рабочее название ТПМ-1 (тенгварская пишущая машинка, модель 1).}
\end{titlepage}

\tableofcontents
\newpage

%=====================================================================
\section{Назначение и исходные положения}

Машинка печатает тенгвар слева направо в режимах, где гласные (а в части режимов и некоторые согласные) передаются надстрочными и подстрочными знаками — техтами. Полные режимы (например, режим Белерианда), где гласные — отдельные тенгвы, поддерживаются автоматически: для них нужны только продвигающие литеры.

\paragraph{Требования.}
\begin{enumerate}[label=Т\arabic*.]
  \item Техта печатается на той же горизонтальной позиции, что и её тенгва, — над лувой (дугой), а не над краем клетки, причём для узких, широких и «смещённых» тенгв.
  \item Межбуквенные промежутки соответствуют ширине тенгв: узкие стоят плотно, широкие получают своё место.
  \item Знаки под тенгвой (черта удвоения, точки снизу) печатаются под лувой, обходя опущенную ножку; черта удвоения при lambe и alda уходит внутрь изгиба.
  \item Порядок ударов не зависит от режима: квенийский (техта над предшествующим согласным) и синдаринский (над последующим) набираются одинаково.
  \item Весь набор литер помещается на обычную клавиатуру из 44 клавиш с двумя регистрами.
\end{enumerate}

\paragraph{Чего машинка не делает.} Она не транскрибирует: наборщик сам знает режим и бьёт тенгвы и техты. Она не кернит: промежутки задаются шириной клетки, а не парами букв. Двенадцатеричные цифры тенгвара пишутся младшим разрядом влево, и наборщик бьёт их в этом порядке — так же, как пишет рука.

%=====================================================================
\section{Письмо глазами механика}

\subsection{Анатомия тенгвы}

Основные тенгвы складываются из ножки (\emph{telco}) и дуги (\emph{lúva}). Ножка бывает висячей, поднятой, короткой или удлинённой; дуга — одинарной или двойной, открытой или закрытой, стоит справа от ножки (столбцы tinco, parma) или слева (calma, quesse). Для машинки существенны три величины:

\begin{itemize}
  \item \textbf{ширина клетки} $w$ — сколько каретке пройти на эту литеру;
  \item \textbf{верхняя ось} — вертикаль, на которой стоит техта; у тенгв с поднятой или опущенной ножкой она не совпадает с серединой клетки;
  \item \textbf{нижняя ось} — вертикаль, на которой стоят подстрочные знаки; она уходит от верхней в сторону, противоположную опущенной ножке.
\end{itemize}

\begin{figure}[h]
\centering
\setlength\tgem{27mm}
\begin{tikzpicture}[x=\tgu,y=\tgu]
\input{gen/fig_anatomy.tex}
\end{tikzpicture}
\caption{Клетки и оси некоторых литер (оттиск, масштаб: 1~em = 27~мм). Серое поле — клетка ширины $w$; штрихпунктир — верхняя ось (техта), пунктир — нижняя ось (подстрочные знаки). Числа — в единицах подачи $u = \mathrm{em}/16$. Знак $d$ — смещение нижней оси от верхней: у tinco ножка слева и нижняя ось уходит вправо, у harma — наоборот.}
\label{fig:anatomy}
\end{figure}

\subsection{Откуда взяты числа}

Все размеры в документе измерены на шрифте \emph{Alcarin Tengwar} (Тоси Омагари, лицензия OFL 1.1, кодировка ConScript Unicode Registry, область U+E000--E07F). В шрифте есть таблица GPOS с якорями: для каждой тенгвы указано, куда ставится верхняя техта и куда — подстрочный знак. Эти якоря и есть «оси» в смысле этого документа. Скрипт \code{tools/measure.py} читает ширины и якоря, \code{tools/gen.py} пересчитывает их в единицы подачи и строит все таблицы и схемы (приложение~\ref{app:method}).

Тот же шрифт подтверждает, что ваши пункты 3 и 4 описывают настоящую необходимость, а не причуду: у черты удвоения в Alcarin пять форм (обычная, укороченная, укороченная со сдвигом, широкая и особая \code{.lambe}, поднятая внутрь изгиба), у носовой черты — шесть, включая особую для silme. В раскладке Дэна Смита, де-факто стандарте восьмибитных тенгварских шрифтов, у каждой гласной техты четыре клавиши — по одной на класс тенгв (широкие, средние, с поднятой ножкой, узкие), а верхняя и нижняя черты даны в узком и широком вариантах. Машинка ТПМ-1 решает ту же задачу иначе: вместо умножения знаков она двигает каретку.

\subsection{Классы ширины}

Если перевести ширины в единицы $u = \mathrm{em}/16$, тенгвы сами собой собираются в четыре группы:

\begin{center}\small
\begin{tabular}{@{}llll@{}}
\toprule
Класс & $w$, u & Кто входит & Сколько \\
\midrule
носитель & 4 & telco, ára, halla & 3 \\
дополнительная & 7--9 & rómen, arda, silme, esse, hyarmen, yanta, úre, sa-rincë & 11 \\
узкая & 10--11 & одинарная дуга (tinco, thúle, óre\ldots), lambe, alda, hwesta sindarinwa & 19 \\
широкая & 15--16 & двойная дуга (ando, anto, númen\ldots) & 16 \\
\bottomrule
\end{tabular}
\end{center}

Отношение широкой к узкой — около 1,6. Это слишком много, чтобы закрыть глаза и сделать машинку моноширинной без заметного уродства (такой вариант всё же разобран в~\S\ref{sec:variants}).

%=====================================================================
\section{Что уже было придумано}

\begin{description}[style=nextline,leftmargin=0pt]
  \item[Тенгварская машинка Кирмсе (1984).] Статья Рональда Кирмсе «QWERTY meets Feanor: The Tengwar Typewriter» в бюллетене \emph{Quettar} №~21. Текста в сети нет; известно, что она вдохновила шрифт Tengwar345. Физических тенгварских машинок я не нашёл.
  \item[Раскладка Дэна Смита (1990-е).] Таблица Приложения~E разложена по рядам клавиатуры; у техт по четыре позиционных варианта, у черт — узкий и широкий. Это «виртуальная машинка» с умножением знаков: наборщик сам выбирает вариант техты под ширину тенгвы.
  \item[Арабская машинка Хаддада (1899).] Четыре формы буквы сведены к двум, потому что соединительная черта целиком отнесена к первой букве, а литеры машинки могут печататься внахлёст. Урок для нас: машинка выигрывает, сокращая число литер за счёт наложения. Ранние арабские машинки огласовок не печатали вовсе.
  \item[Еврейская раскладка.] Огласовки (никуд) вынесены на третий слой клавиатуры, отдельно от букв.
  \item[IBM Executive (1944).] Пропорциональная подача: литеры занимают от двух до пяти единиц по 1/32 дюйма, пробел разделён на клавиши в 2 и 3 единицы. Это прямой предок механизма подачи ТПМ-1.
  \item[Ходовой винт (IBM, 1960-е).] Каретку двигает винт, поворот которого пропорционален ширине литеры. Такой привод снимает ограничение на мелкость шага, которое есть у зубчатой рейки.
\end{description}

%=====================================================================
\section{Принцип: осевые литеры}

\subsection{Почему не годится классическая мёртвая клавиша}

На обычной машинке литера отлита так, что её левый край — у края клетки. Мёртвая клавиша (акцент) не двигает каретку и бьёт в фиксированное место клетки, рассчитанное на «среднюю» букву. Для тенгвара этого мало: середина узкой клетки не совпадает ни с серединой широкой, ни с осью harma, ни тем более с осью hyarmen. Отсюда и четыре набора техт у Смита.

\begin{figure}[h]
\centering
\setlength\tgem{19mm}
\begin{tikzpicture}[x=\tgu,y=\tgu]
\input{gen/fig_principle.tex}
\end{tikzpicture}
\caption{Одна и та же клавиша «e» на tinco, ando, harma, hyarmen. Вверху — классика: литера у левого края клетки, мёртвая клавиша бьёт в 5\,u от края (середина узкой клетки). Красным — промахи. Внизу — осевые литеры: каретка ставит под удар ось буквы, и одна и та же техта садится верно везде.}
\label{fig:principle}
\end{figure}

\subsection{Осевая литера}

\textbf{Правило отливки.} Каждая продвигающая литера отлита так, что её верхняя ось проходит через ось рычага (точку удара). Каждая техта отлита так, что её собственная ось (якорь знака) проходит через ту же точку. Тогда техта всегда садится на ось буквы, где бы эта ось ни была внутри клетки.

\textbf{Правило подачи.} Клетка литеры делится осью на две части: $a$ — от левого края до оси, $b$ — от оси до правого края, $a + b = w$. Каретка проходит $a$ до удара литеры и $b$ — после. Мёртвые знаки ($w = 0$) бьют в текущую точку и каретку не двигают.

\subsection{Отложенная вторая фаза}

Если выполнять $b$ сразу после удара, каретка уйдёт с оси раньше, чем наборщик ударит техту. Поэтому вторая фаза \emph{откладывается}: величина $b$ запоминается в регистре и выполняется в начале следующей продвигающей литеры (или пробела) вместе с её первой фазой. Каретка проходит за один цикл
\[
  \Delta = b_{\text{предыдущей}} + a_{\text{текущей}}
\]
и стоит на оси текущей буквы до тех пор, пока не придёт следующая продвигающая клавиша. Все техты, удвоения и прочие знаки, набранные в промежутке, садятся на эту ось.

Отсюда требование Т4: наборщик \emph{всегда} бьёт сначала тенгву, потом её знаки, независимо от того, к какому звуку знак относится. В квенийском режиме это совпадает с порядком чтения; в синдаринском (гласная над следующим согласным) наборщик просто бьёт согласный раньше, чем гласную над ним, как это и делает рука.

\subsection{Нижняя поправка}

Подстрочные знаки стоят на нижней оси, а она у тенгв с опущенной ножкой смещена от верхней на одну единицу и больше (рис.~\ref{fig:anatomy}). В коде каждой тенгвы есть поправка $d$ — целое число единиц от верхней оси до нижней. Удар подстрочного знака выполняется так: каретка сдвигается на $d$, рычаг бьёт, каретка возвращается на $-d$. Регистр при этом не меняется, так что после подстрочного знака можно бить надстрочный, и наоборот.

Распределение $d$ по 48 тенгвам: $d=0$ — 22, $\pm 1$ — 20, $\pm 2$ — 5, $-3$ — 1 (hyarmen). Знак $d$ прямо показывает сторону опущенной ножки: у столбцов tinco и parma он положительный, у calma и quesse — отрицательный.

%=====================================================================
\section{Единицы, размеры, точность}

\begin{center}\small
\begin{tabular}{@{}lll@{}}
\toprule
Величина & Значение & Замечание \\
\midrule
кегль (em) & 1/4 дюйма = 6,35 мм & высота дуги $\approx$ 2,6 мм, как строчные у пики \\
единица подачи $u$ & em/16 = 1/64 дюйма $\approx$ 0,397 мм & вдвое мельче IBM Executive \\
ширины & 4--16\,u (1,6--6,4 мм) & широкая литера крупнее обычной \\
$a$, $b$ & 0--9\,u & по 4 бита \\
$d$ & $-3\ldots+2$\,u & 3 бита со знаком \\
пробелы & 4\,u и 6\,u & две клавиши, как у Executive \\
межстрочие & 1/3 дюйма $\approx$ 8,47 мм & «двойной» интервал: литеры с ножками и техтами занимают $\approx$ 1,05 em \\
\bottomrule
\end{tabular}
\end{center}

\paragraph{Погрешности квантования.} Литера отлита по точному якорю, поэтому \emph{техта стоит на оси точно}: квантование задевает только края клетки. У тенгв край клетки уходит от края знака не больше чем на 0,5\,u ($\approx$ 0,2 мм) — это и есть погрешность межбуквенного промежутка. Нижняя поправка $d$ округлена до целой единицы; остаточная ошибка подстрочного знака не превышает 0,4\,u ($\approx$ 0,16 мм). Обе цифры — порядка допуска выравнивания рычагов на обычной машинке.

\paragraph{Вертикаль.} Почти у всех тенгв техта стоит на одной высоте (якорь на 0,40 em над строкой). Исключения — silme и esse (0,52 em) и halla (0,70 em). Над silme и esse техту и в письме не ставят — для этого существуют формы \emph{nuquerna}, а halla гласных не несёт. Поэтому все надстрочные литеры отлиты на одной высоте, а поворот валика на полстроки, если он понадобится, — штатная функция любой машинки.

%=====================================================================
\section{Общая схема машинки}

ТПМ-1 — электрическая рычажная машинка с силовым валом (как IBM Electric) и неподвижной точкой удара. Её отличие от обычной — кодовый тракт подачи: каждая клавиша сообщает не просто «сдвинь на клетку», а трёхчастный код $(a, b, d)$, и каретку двигает сумматор с регистром.

\begin{figure}[h]
\centering
\begin{tikzpicture}[node distance=7mm and 9mm]
\node[blk] (key) {Клавиша\\[-1pt]\footnotesize рычаг клавиши};
\node[blk,right=of key] (comb) {Кодовая гребёнка\\[-1pt]\footnotesize 2 яруса $\times$ 12 выступов};
\node[blk,right=of comb] (bails) {Считывающие планки\\[-1pt]\footnotesize 12 шт., рамка с подъёмом};
\node[blk,below=12mm of bails] (latch) {Защёлки кода\\[-1pt]\footnotesize $A,\ B,\ D,\ N$};
\node[blk,left=of latch] (reg) {Регистр R\\[-1pt]\footnotesize $a_{\text{посл}},\ b,\ d$};
\node[blk,left=of reg] (sum) {Сумматор\\[-1pt]\footnotesize $\Delta = R.b + A$\\[-1pt]\footnotesize или $\pm R.d$};
\node[blk,below=12mm of sum] (car) {Привод каретки\\[-1pt]\footnotesize ходовой винт, шаг 1/64''};
\node[blk,below=12mm of latch] (bar) {Литерный рычаг\\[-1pt]\footnotesize удар в ось};
\node[blk,below=12mm of reg] (cyc) {Силовой вал\\[-1pt]\footnotesize цикл 360°};
\node[blk,above=8mm of bails] (shift) {Регистр (Shift)\\[-1pt]\footnotesize сегмент + рамка планок};
\draw[arr] (key) -- (comb);
\draw[arr] (comb) -- (bails);
\draw[arr] (shift) -- (bails);
\draw[arr] (bails) -- (latch);
\draw[arr] (latch) -- (reg);
\draw[arr] (reg) -- (sum);
\draw[arr] (latch.south west) to[bend left=10] (sum.south east);
\draw[arr] (sum) -- (car);
\draw[arr] (key.south) |- (cyc.west);
\draw[arr] (cyc) -- (bar);
\draw[arr,dashed] (cyc) -- node[lbl,right,pos=.45]{такт} (reg);
\node[lbl,below=1mm of car,text width=40mm,align=center] {перемещение $\Delta$ — до удара};
\node[lbl,below=1mm of bar,text width=40mm,align=center] {удар после остановки каретки};
\end{tikzpicture}
\caption{Кодовый тракт ТПМ-1. Сплошные стрелки — передача кода и движения, пунктир — тактирование от силового вала.}
\label{fig:blocks}
\end{figure}

\begin{figure}[h]
\centering
\begin{tikzpicture}[x=1mm,y=1mm,line cap=round]
% base
\draw[mech,fill=gray!8] (0,0) rectangle (150,4);
% keylever
\draw[mech] (8,14) -- (88,22);
\fill (88,22) circle (1.1);
\node[lbl,anchor=south] at (88,23) {ось рычага};
\draw[mech,fill=white] (4,12) rectangle (14,17); \node[lbl] at (9,9) {клавиша};
% code comb
\foreach \i in {0,...,5} {\draw[mech,fill=dead] (30+\i*5,15.4+\i*0.5) rectangle +(3,{-3-mod(\i,2)*2});}
\draw[decorate,decoration={brace,amplitude=4pt}] (29,8) -- (59,8) node[midway,lbl,below=3pt]{кодовая гребёнка (выступы $A$, $B$, $D$, $N$)};
% bails
\foreach \i in {0,...,5} {\draw[mech,fill=adv] (31.5+\i*5,3+\i*0.0) ++(0,2) circle (1.3);}
\node[lbl,anchor=west] at (61,6.5) {считывающие планки};
% trip / power roll
\draw[mech,fill=gray!15] (100,10) circle (5); \node[lbl] at (100,10) {вал};
\draw[mech] (88,22) -- (97,14);
% typebar
\draw[mech] (112,12) -- (128,52);
\fill (112,12) circle (1.1);
\draw[mech,fill=ink!80] (126.5,50) rectangle (130.5,55);
\node[lbl,anchor=west] at (131,52) {литера};
\draw[mech] (97,14) to[out=20,in=-120] (112,12);
% platen and carriage
\draw[mech,fill=gray!20] (133,64) circle (8);
\node[lbl] at (133,64) {валик};
\draw[mech,fill=gray!6] (110,74) rectangle (156,80); \node[lbl] at (133,77) {каретка};
\draw[mech] (100,84) -- (160,84);
\foreach \x in {100,102,...,160} {\draw[ink!60] (\x,83) -- (\x+1,85);}
\node[lbl,anchor=south] at (130,85.5) {ходовой винт};
\fill[axisred] (131,57) circle (0.9); \node[lbl,text=axisred,anchor=east] at (124,58) {точка удара};
% register box
\draw[mech,fill=white,rounded corners=1pt] (64,40) rectangle (96,56);
\node[lbl,align=center] at (80,48) {защёлки + регистр R\\сумматор $\Delta$};
\draw[arr] (55,14) -- (68,40);
\draw[arr] (96,52) -- (104,80);
\end{tikzpicture}
\caption{Общая схема, вид сбоку (условно, без масштаба). Рычаг клавиши несёт кодовую гребёнку; при нажатии он расцепляет силовой вал, вал поднимает литерный рычаг; одновременно выступы гребёнки толкают считывающие планки, код запоминается в защёлках, и сумматор поворачивает ходовой винт каретки ещё до удара.}
\label{fig:side}
\end{figure}

\subsection{Кодовое слово}

Каждый ярус кодовой гребёнки несёт 12 позиций выступов:

\begin{center}\small
\begin{tabular}{@{}lccl@{}}
\toprule
Поле & Биты & Диапазон & Смысл \\
\midrule
$A$ & 4 & 0--15 & первая фаза: от края клетки до оси \\
$B$ & 4 & 0--15 & вторая (отложенная) фаза: от оси до края \\
$D$ & 3 & $-4\ldots+3$ & нижняя поправка (запоминается вместе с $B$) \\
$N$ & 1 & 0/1 & подстрочный знак: бить со сдвигом на $R.d$ \\
\bottomrule
\end{tabular}
\end{center}

Продвигающая литера: $A + B > 0$. Надстрочная техта: всё по нулям. Подстрочный знак: $N = 1$, остальное — нули. Пробелы и служебные клавиши кода не несут: их действие жёстко задано.

\subsection{Два яруса гребёнки}

На обычной рычажной машинке строчная и прописная литеры сидят на одном рычаге, и их ширина обязана совпадать (или подача от неё не зависит). В ТПМ-1 так не выйдет: на одной клавише окажутся, например, гласная техта ($w=0$) и ára с той же техтой ($w=4$). Поэтому у гребёнки два яруса выступов, а клавиша регистра поднимает не только сегмент рычагов, но и рамку считывающих планок, так что планки читают верхний ярус. Это избавляет раскладку от требования «одинаковый код на одной клавише» и позволяет ставить на регистр то, что связано по смыслу (\S\ref{sec:keyboard}).

\begin{figure}[h]
\centering
\begin{tikzpicture}[x=1mm,y=1mm]
\foreach \t/\yy/\nm in {0/0/нижний ярус,1/10/верхний ярус} {
  \node[lbl,anchor=east] at (-3,\yy+2.5) {\nm};
  \foreach \i in {0,...,11} {\draw[ink!50] (\i*6,\yy) rectangle +(5,5);}
}
% example code: tinco lower (A=5,B=5,D=+1), ext tinco upper (A=6,B=4,D=0)
\def\bits#1#2{\foreach \i in {#1} {\fill[ink!80] (\i*6+0.6,#2+0.6) rectangle +(3.8,3.8);}}
\bits{0,2,4,6,11}{0}
\bits{1,2,6,8}{10}
\foreach \i/\l in {0/1,1/2,2/4,3/8} {\node[lbl] at (\i*6+2.5,-3) {$A_{\l}$};}
\foreach \i/\l in {4/1,5/2,6/4,7/8} {\node[lbl] at (\i*6+2.5,-3) {$B_{\l}$};}
\foreach \i/\l in {8/1,9/2,10/\pm} {\node[lbl] at (\i*6+2.5,-3) {$D_{\l}$};}
\node[lbl] at (11*6+2.5,-3) {$N$};
\node[lbl,anchor=west] at (74,2.5) {tinco: $A{=}5$, $B{=}5$, $D{=}{+}1$};
\node[lbl,anchor=west] at (74,12.5) {удл.~tinco: $A{=}6$, $B{=}4$, $D{=}0$};
\end{tikzpicture}
\caption{Два яруса кодовой гребёнки клавиши Q (tinco / удлинённая tinco). Чёрный квадрат — выступ. Поле $D$ — величина (два бита) и знак. На нижнем ярусе в позиции $N$ стоит маркер вместо знака $D$ — в примере $N$ занят под знак, см. текст.}
\label{fig:comb}
\end{figure}

\noindent\emph{Уточнение к рис.~\ref{fig:comb}.} В примере $D$ кодируется как «величина + знак» (биты $D_1$, $D_2$, $D_\pm$), а позиция $N$ у литер всегда пуста; у tinco выступ в одиннадцатой позиции — это $D_\pm$ = «плюс». Кодировка $D$ (дополнительный код или знак и величина) — дело конструктора сумматора, на раскладку она не влияет.

%=====================================================================
\section{Механизм подачи}

\subsection{Автомат состояний}

Каретка в каждый момент находится в одном из двух состояний: \textbf{на краю} (регистр пуст, точка удара — граница клеток) или \textbf{на оси} (в регистре $a,b,d$ последней литеры, точка удара — её ось).

\begin{figure}[h]
\centering
\begin{tikzpicture}[->,>={Latex[length=2.2mm]},node distance=62mm,thick,
  every state/.style={draw=ink!80,fill=adv,minimum size=22mm,font=\sffamily\small,align=center},
  lab/.style={font=\sffamily\scriptsize,align=center,fill=white,inner sep=1.5pt}]
\node[state,initial,initial text={\scriptsize ВК}] (E) {НА КРАЮ\\\scriptsize $R=0$};
\node[state,right=of E] (A) {НА ОСИ\\\scriptsize $R=(a,b,d)$};
\path (E) edge[bend left=22] node[lab,above]{тенгва: $+a$} (A)
      (A) edge[bend left=22] node[lab,below]{пробел: $+b+s$\\«назад на литеру»: $-a$} (E)
      (A) edge[loop above,looseness=6,out=60,in=120] node[lab,above]{тенгва: $+b+a'$, $R\leftarrow(a',b',d')$} (A)
      (A) edge[loop right,looseness=5,out=-20,in=20] node[lab,right]{техта: 0\\подстр.: $+d$, удар, $-d$\\$\pm u$: $\pm1$, $b\mp1$} (A)
      (E) edge[loop above,looseness=6,out=60,in=120] node[lab,above]{пробел: $+s$; $\pm u$: $\pm1$} (E)
      (E) edge[loop left,looseness=5,out=160,in=200] node[lab,left]{знак без буквы:\\удар на краю} (E);
\end{tikzpicture}
\caption{Автомат каретки. ВК — возврат каретки (обнуляет регистр). $s$ — ширина пробела (4 или 6\,u).}
\label{fig:fsm}
\end{figure}

\paragraph{Служебные клавиши.}
\begin{itemize}
  \item \textbf{Пробел 4u и 6u.} Сначала выполняют отложенную фазу $b$, затем свой сдвиг; регистр обнуляется.
  \item \textbf{«Назад на литеру».} Отводит каретку на $a$ последней литеры, к левому краю её клетки, и обнуляет $b$. Если затем ударить ту же литеру (например, через корректирующую ленту), она вернётся точно на ось. Глубина памяти — одна литера: для более дальнего возврата есть только шаги по 1\,u.
  \item \textbf{$\leftarrow$u и u$\rightarrow$.} Сдвигают каретку на одну единицу и одновременно поправляют $b$ в регистре на обратную величину, так что следующая буква встанет как обычно. Это ручная поправка оси для редких случаев (техта над нетипичной формой, подстрочный знак под alda, разрядка) — та самая клавиша-модификатор из исходной постановки, но нужная лишь изредка.
\end{itemize}

\subsection{Циклограмма}

Силовой вал делает один оборот на каждое нажатие. Перемещение каретки целиком укладывается в первую половину оборота, удар — после остановки каретки.

\begin{figure}[h]
\centering
\begin{tikzpicture}[x=0.36mm,y=1mm]
\draw[ink!50,->] (0,0) -- (380,0) node[lbl,right]{угол вала};
\foreach \a in {0,90,180,270,360} {\draw[ink!50] (\a,0) -- (\a,-1.5) node[lbl,below]{\a°};}
\def\row#1#2{\node[lbl,anchor=east] at (-4,#1) {#2};}
\row{38}{чтение кода}
\fill[accent!35] (10,36) rectangle (70,40);
\row{30}{подача $\Delta$ (литера)}
\fill[axisblue!40] (80,28) rectangle (190,32);
\row{22}{сдвиг $+d$ (подстр. знак)}
\fill[axisblue!25] (80,20) rectangle (150,24);
\row{14}{удар рычага}
\fill[axisred!55] (205,12) rectangle (245,16);
\row{6}{возврат $-d$; запись R}
\fill[axisblue!25] (255,4) rectangle (320,8);
\fill[accent!20] (255,8) rectangle (340,9.5);
\draw[ink!40,dashed] (200,0) -- (200,42);
\node[lbl,anchor=south] at (200,42) {каретка неподвижна};
\end{tikzpicture}
\caption{Циклограмма одного оборота силового вала. Для продвигающей литеры работают строки 1, 2, 4 и запись регистра; для подстрочного знака — 1, 3, 4 и возврат; для надстрочной техты — только 1 и 4.}
\label{fig:timing}
\end{figure}

\paragraph{Быстрый набор.} Рычаги машинки с силовым валом срабатывают по очереди, даже если клавиши перекрываются, а отложенная фаза выполняется в начале \emph{следующего} цикла. Поэтому перекрытие нажатий («перекат») не путает оси. Единственное правило для наборщика — знаки бьются после своей тенгвы.

%=====================================================================
\section{Типы литер}
\label{sec:types}

Литеры делятся на семь типов по поведению каретки и способу отливки. На рисунках — \emph{оттиск} (лицо литеры зеркально ему), рамка — контур литеры, серое поле — клетка, красный штрихпунктир — ось рычага, синий пунктир — нижняя ось ($a + d$).

\subsection*{Тип I. Продвигающая тенгва}

Код $(a, b, d)$, $N = 0$. Знак отлит по верхнему якорю, края клетки округлены до единицы. Сюда же относятся sa-rincë (концевой крюк ставится отдельной продвигающей литерой после буквы, а не надстрочным знаком), цифры и знаки препинания (тип~VII).

\begin{figure}[h]
\centering
\setlength\tgem{30mm}
\begin{tikzpicture}[x=\tgu,y=\tgu]
\input{gen/fig_slugs_T.tex}
\end{tikzpicture}
\caption{Литеры типа I: tinco ($5+5$), ando ($8+8$), harma ($4+6$, ось левее середины), hyarmen ($7+2$, ось почти у правого края; $d = -3$), короткий носитель ($2+2$).}
\end{figure}

\subsection*{Тип II. Надстрочная мёртвая техта}

Гласные a, e, i, o, u и тильда. $w = 0$, отлита по якорю знака. Один комплект на все тенгвы — это главная экономия ТПМ-1 по сравнению с раскладкой Смита. Высота одна для всех (\S5).

\subsection*{Тип III. Подстрочный знак постоянной длины}

Точка снизу, две точки снизу (y), метка младшего разряда цифр. $N = 1$: бьются со сдвигом на $R.d$.

\subsection*{Тип IV. Мерный знак}

Черта удвоения и носовая черта — в узком (7,2\,u) и широком (12,6\,u) вариантах. Длина зависит от числа дуг, поэтому центрирование её не решает, и два варианта остаются — ваш пункт~4 в чистом виде. Укороченные формы шрифта (под носитель) машинке не нужны: удвоенный носитель не встречается.

\subsection*{Тип V. Особый знак: удвоение для lambe и alda}

Черта, отлитая выше строки (0,17--0,23~em), чтобы войти внутрь нижнего изгиба lambe. Подстрочный ($N=1$), но поднятый. У lambe нижняя ось совпадает с верхней ($d=0$), у alda смещена на +1\,u.

\begin{figure}[h]
\centering
\setlength\tgem{24mm}
\begin{tikzpicture}[x=\tgu,y=\tgu]
\input{gen/fig_slugs_D.tex}
\end{tikzpicture}
\caption{Мёртвые литеры типов II--V: техты a, e, o; две точки снизу; черта удвоения узкая, широкая и особая для lambe. Ширина клетки нулевая, рамка показывает только размер литеры.}
\end{figure}

\subsection*{Тип VI. Комбинированная литера ára + техта}

Длинный носитель с уже отлитой гласной: долгие á, é, í, ó, ú одним ударом. Код $(2, 2, 0)$. Стоят на регистре гласных клавиш: Shift + «e» = долгое é на ára.

\subsection*{Тип VII. Цифры и знаки препинания}

Продвигающие, код по якорю, если он есть в шрифте (у цифр он есть — над цифрой ставится черта или точка), иначе по середине. Ширины цифр в Alcarin заметно различаются (5--16\,u); для машинного шрифта их стоит выровнять (\S\ref{sec:risks}).

\begin{figure}[h]
\centering
\setlength\tgem{30mm}
\begin{tikzpicture}[x=\tgu,y=\tgu]
\input{gen/fig_slugs_C.tex}
\end{tikzpicture}
\caption{Литеры типов VI и VII и концевой крюк sa-rincë.}
\end{figure}

%=====================================================================
\section{Раскладка клавиатуры}
\label{sec:keyboard}

44 клавиши в четыре ряда по 11, два регистра — 88 мест; занято 86, два в резерве. Логика:

\begin{itemize}
  \item \textbf{Ряд цифр}: двенадцатеричные цифры на нижнем регистре, знаки препинания и цифра 11 — на верхнем.
  \item \textbf{Верхний буквенный ряд}: тенгвы с висячей ножкой (tinco, parma, calma, quesse, ando, umbar), anca, anto, ampa, rómen, lambe; на регистре — их удлинённые формы (и anga, arda, alda).
  \item \textbf{Основной ряд}: слева согласные (thúle, formen, harma, númen, malta), справа — гласные техты. Слог тенгварского письма — «тенгва + техта», так что руки чередуются, как при слепом наборе. Shift + гласная — та же гласная на ára.
  \item \textbf{Нижний ряд}: óre, vala, anna, silme и silme nuquerna, hyarmen, yanta, оба носителя и знаки удвоения и носовой черты.
  \item \textbf{Служебный ряд}: регистр, два пробела, $\leftarrow$u, u$\rightarrow$, «назад на литеру». Возврат каретки и табулятор — как обычно (на схеме не показаны).
\end{itemize}

\begin{figure}[h]
\centering
\begin{tikzpicture}[x=1mm,y=1mm]
\input{gen/keyboard.tex}
\end{tikzpicture}
\caption{Раскладка ТПМ-1. Внизу клавиши — нижний регистр, вверху фиолетовым — верхний. Мёртвые знаки показаны на пунктирной оси. Мелкие латинские буквы — соответствие позициям QWERTY, только для ориентировки.}
\label{fig:kbd}
\end{figure}

Раскладка составлена по структуре письма, а не по частотам. Если собрать корпус квенийских и синдаринских текстов, перестановка внутри рядов по частоте — отдельная небольшая задача; кодовые гребёнки от неё не меняются, меняются только надписи и литеры.

%=====================================================================
\section{Набор: два примера}

Ниже — пословный разбор двух строк. Таблицы и оттиски построены симулятором каретки (\code{simulate} в \code{tools/gen.py}) по кодам литер; эталон под оттиском набран тем же шрифтом с полной OpenType-раскладкой. Совпадение посадки техт и подстрочных знаков проверяет весь принцип; различия в промежутках — это квантование до 1\,u и отсутствие кернинга.

\subsection{Квенья: \emph{Elen síla lúmenn' omentielvo}}

\begin{figure}[h]
\centering
\setlength\tgem{12.5mm}
\begin{tikzpicture}[x=\tgu,y=\tgu]
\input{gen/fig_trace_q.tex}
\end{tikzpicture}
\caption{Квенийский режим: техта над предшествующим согласным. Короткие красные риски под строкой — точки удара тенгв.}
\end{figure}

{\small\setlength\tgem{5mm}
\begin{longtable}{@{}rlrrl@{}}
\toprule
\# & Удар & Подача, u & Ось, u & Регистр $b, d, a$ \\
\midrule\endhead
\input{gen/tab_trace_q.tex}
\bottomrule
\end{longtable}}

\subsection{Синдарин: \emph{pedo mellon a minno}}

\begin{figure}[h]
\centering
\setlength\tgem{14mm}
\begin{tikzpicture}[x=\tgu,y=\tgu]
\input{gen/fig_trace_s.tex}
\end{tikzpicture}
\caption{Синдаринский режим (общий): техта над последующим согласным. В «mellon» под lambe — особая черта удвоения (тип V), в «minno» под númen — узкая (тип IV).}
\end{figure}

{\small\setlength\tgem{5mm}
\begin{longtable}{@{}rlrrl@{}}
\toprule
\# & Удар & Подача, u & Ось, u & Регистр $b, d, a$ \\
\midrule\endhead
\input{gen/tab_trace_s.tex}
\bottomrule
\end{longtable}}

\noindent Видно, что порядок ударов в обоих режимах одинаков — «тенгва, затем её знаки», хотя в квенийском «e» в \emph{elen} читается после lambe, а в синдаринском «e» в \emph{mellon} — перед ней.

%=====================================================================
\section{Варианты исполнения}
\label{sec:variants}

\subsection*{Вариант Б. Моноширинный}

Все клетки одной ширины (например, 12\,u), литеры перерисованы: широкие сжаты на четверть, узкие разрежены, носители получают поля. Ось — посередине клетки, подача фиксированная, а отложенная вторая фаза сводится к постоянному полушагу. Механика — почти серийная. Цена: сжатые двойные дуги, разрежённые носители, и строка выходит примерно на 15\% длиннее, чем при пропорциональной подаче. Хороший вариант для первого макета.

\subsection*{Вариант В. Составной (по Хаддаду)}

32 основные и удлинённые тенгвы — это 4 вида ножки $\times$ 4 столбца $\times$ одинарная или двойная дуга. Можно отлить отдельно ножки (мёртвые) и дуги (продвигающие) и печатать букву двумя ударами. 12 клавиш вместо 32 освобождают 20 мест — например, под второй набор техт для свободной вёрстки. Цена: вдвое больше ударов на букву, требования к точности стыка ножки и дуги, а ножка стоит слева у столбцов tinco и parma и справа у calma и quesse (это видно по знаку $d$), так что литер-ножек нужно восемь, а не четыре. Интересен как демонстрация того, что письмо Феанора «собрано» из деталей, но для повседневного набора хуже основного варианта.

\subsection*{Вариант Г. Классический (по Смиту)}

Обычная подача по левому краю, четыре набора техт и узкие/широкие черты, как в раскладке Смита. Механика серийная, но знаков становится больше, чем мест: гласных $5\times4 = 20$ вместо 5, плюс узкие и широкие черты. На 88 мест это уже не помещается, нужен третий регистр. Плюс наборщик сам выбирает вариант техты.

\begin{center}\small
\begin{tabular}{@{}lcccc@{}}
\toprule
& А (основной) & Б (моно) & В (составной) & Г (по Смиту) \\
\midrule
техт гласных & 5 & 5 & 5 & 20 \\
клавиш на 32 тенгвы & 32 & 32 & 12 & 32 \\
ударов на слог «тенгва+техта» & 2 & 2 & 3 & 2 \\
выбор варианта знака & нет & нет & нет & вручную \\
механика подачи & новая & почти серийная & новая & серийная \\
верность пропорциям & высокая & низкая & высокая & высокая \\
\bottomrule
\end{tabular}
\end{center}

%=====================================================================
\section{Риски и открытые вопросы}
\label{sec:risks}

\begin{enumerate}[label=Р\arabic*.]
  \item \textbf{Мелкий шаг.} 1/64 дюйма вдвое мельче шага IBM Executive. Зубчатая рейка с таким шагом ненадёжна; отсюда привод ходовым винтом. Альтернатива — крупнее кегль (em = 1/2 дюйма), но тогда литеры становятся плакатными.
  \item \textbf{Двухъярусная гребёнка.} Подъём рамки считывающих планок вместе с сегментом — нештатное решение. Если оно не пройдёт, можно вернуться к правилу «одинаковый код на клавише» и переставить пары: для большинства тенгв подходящие пары есть, хуже всего с парами «гласная / ára+гласная».
  \item \textbf{Масса широкой литеры.} 6,4 мм лица против $\approx$2,5 мм у пики: тяжелее рычаг, сильнее удар. Нужна регулировка силы удара по рычагам (на электрических машинках она обычно есть).
  \item \textbf{Память на одну литеру.} «Назад на литеру» работает только на одну букву. Для исправлений глубже — шаги по 1\,u и глазомер.
  \item \textbf{Подстрочные знаки под alda и hyarmen.} У alda два разных нижних якоря (для черты и для точек), а $d$ один; точки под alda встанут с ошибкой около 1,2\,u. У hyarmen $d = -3$. Оба случая редки; есть клавиши $\pm u$.
  \item \textbf{o и u.} Принятая здесь пара «правый завиток = o, левый = u» в ряде режимов меняется местами. На машинке это просто надписи на клавишах.
  \item \textbf{Цифры.} В Alcarin ширины цифр разные; для машинки лучше выровнять их по двум-трём значениям. Порядок разрядов (младший слева) наборщик соблюдает сам.
  \item \textbf{Машинный рисунок.} Alcarin — книжный шрифт с контрастом. Для литер нужен свой рисунок: равномерный штрих, утолщённые техты, укрупнённые точки. Метрики (клетки, оси) при этом можно оставить.
\end{enumerate}

\paragraph{Нужно решить.} Кегль и шаг (Р1); нужны ли удлинённые тенгвы в основном наборе (они занимают 8 мест верхнего регистра); оставлять ли комбинированные ára-литеры; собирать ли частотный корпус для раскладки.

%=====================================================================
\appendix
\section{Полная таблица литер}
\label{app:tables}

$w, a, b, d$ — в единицах $u = \mathrm{em}/16$; «шир.» — ширина знака в шрифте (единицы em/1000); «края» — на сколько край клетки слева/справа отстоит от края знака (плюс — клетка шире знака); «ост.~$d$» — остаточная ошибка нижней оси после округления.

{\small\setlength\tgem{5.5mm}
\begin{longtable}{@{}cllrrrrrcr@{}}
\toprule
& Тенгва & U+ & шир. & $w$ & $a$ & $b$ & $d$ & края & ост.~$d$ \\
\midrule\endhead
\input{gen/tab_tengwar.tex}
\bottomrule
\end{longtable}}

{\small\setlength\tgem{5.5mm}
\begin{longtable}{@{}cllrrrr@{}}
\toprule
& Литера & U+ & шир. & $w$ & $a$ & $b$ \\
\midrule\endhead
\input{gen/tab_others.tex}
\bottomrule
\end{longtable}}

{\small\setlength\tgem{5.5mm}
\begin{longtable}{@{}cllcrc@{}}
\toprule
& Знак & Глиф шрифта & Положение & Длина, u & $N$ \\
\midrule\endhead
\input{gen/tab_marks.tex}
\bottomrule
\end{longtable}}

\section{Методика измерений}
\label{app:method}

\begin{enumerate}
  \item Ширина клетки — продвижение глифа (\code{hmtx}); $w = \mathrm{round}(\text{adv}/62{,}5)$.
  \item Верхняя ось — базовый якорь класса «над» в поиске GPOS MarkToBase (поиск 6 в Alcarin 0.83); $a = \mathrm{round}(x_{\text{верх}}/62{,}5)$, $b = w - a$.
  \item Нижняя ось — базовый якорь класса «под» в том же поиске; $d = \mathrm{round}((x_{\text{низ}} - x_{\text{верх}})/62{,}5)$. Для lambe и alda дополнительно читается отдельный поиск с вариантами \code{.lambe}.
  \item Ось знака — якорь знака (MarkRecord); литера техты отливается так, чтобы этот якорь лёг на ось рычага.
  \item Симулятор: регистр $R=(a,b,d)$, правила рис.~\ref{fig:fsm}; оттиск строится установкой глифов по рассчитанным осям, эталон — набором Unicode-строки движком HarfBuzz.
\end{enumerate}

Файлы: \code{tools/inventory.py} (набор литер), \code{tools/layout.py} (раскладка), \code{tools/measure.py} (замеры $\to$ \code{data/codes.json}), \code{tools/gen.py} (таблицы, схемы, симуляция $\to$ \code{gen/}), \code{main.tex}. Пересборка: \code{make}.

\section*{Источники}
\begin{itemize}\small
  \item R.\,E.~Kyrmse. QWERTY meets Feanor: The Tengwar Typewriter // Quettar, №~21, 1984, с.~3--5 (по ссылке: \url{https://sites.usp.br/otto/?p=182}).
  \item D.\,S.~Smith. Tengwar Quenya — документация шрифта (раскладка, позиции техт).
  \item Tolkien Gateway: Dan Smith — \url{https://tolkiengateway.net/wiki/Dan_Smith}.
  \item T.~Omagari. Alcarin Tengwar — \url{https://github.com/Tosche/Alcarin-Tengwar} (OFL 1.1).
  \item Amanye Tenceli: The Tengwar — \url{https://at.boktypografen.se/tengwar.htm}.
  \item Tecendil — \url{https://www.tecendil.com/blog}.
  \item J.~Savard. The Invention of the Arabic Typewriter — \url{https://typedrawers.com/discussion/comment/26079}.
  \item Hebrew keyboard — \url{https://en.wikipedia.org/wiki/Hebrew_keyboard}.
  \item IBM Electric — \url{https://en.wikipedia.org/wiki/IBM_Electric}; IBM Model 04 Executive — \url{https://typewriterdatabase.com/1950-ibm-model-04-executive.15625.typewriter}.
  \item US~3346088. Carrier mechanism for proportional escapement typewriter — \url{https://patents.google.com/patent/US3346088A}.
\end{itemize}

\end{document}
